\subsection{תוצאות כמותיות}

הערכת המערכת בוצעה בכמה רמות: זיהוי הרובוטים, זיהוי הספרות, והשיוך
הסופי של עקבה לקבוצה.

כדי להעריך את ביצועי המערכת המלאה מקצה לקצה, בוצע ניסוי אמפירי המבוסס
על מתודולוגיית תיקוף ידני. במסגרת הבדיקה נדגמו מופעי רובוטים במרווחי
זמן קבועים מתוך שלושה מקצי תחרות מלאים משלוש תחרויות שונות, כדי לייצר
``\textenglish{Ground Truth}'' להשוואה. ``מופע רובוט'' הוגדר כרגע שבו
רובוט נראה על המסך במידה שמאפשרת לאדם לעקוב אחריו לאורך הסרטון, גם אם
מספר הקבוצה אינו קריא באותו פריים ספציפי. מופע נספר כהצלחה רק אם
המערכת שייכה באותו רגע את העקבה לקבוצה הנכונה.

מדד הדיוק המערכתי מוגדר כיחס בין מספר המופעים שבהם המערכת שייכה
בהצלחה את הרובוט לקבוצה הנכונה לבין סך המופעים הכולל שנדגם. בניסוי זה
נדגמו בסך הכל 273 מופעי רובוט לאורך שלושת המקצים. המערכת זיהתה ושייכה
בהצלחה ובאופן מדויק את הרובוט לקבוצתו ב־204 מופעים, נתון המוביל למדד
דיוק מערכתי של $\frac{204}{273} \approx 0.75$.

\begin{table}[H]
  \centering
  \begin{tabular}{|c|c|c|c|c|}
    \hline
    \textbf{מקצה} &
    \textbf{מופעי רובוט שנדגמו} &
    \textbf{שיוכים נכונים} &
    \textbf{שיוכים חסרים או שגויים} &
    \textbf{דיוק שיוך} \\ \hline

    1 & 91 & 71 & 20 & 78.0\% \\ \hline
    2 & 88 & 61 & 27 & 69.3\% \\ \hline
    3 & 94 & 72 & 22 & 76.6\% \\ \hline
    \textbf{סה''כ} & \textbf{273} & \textbf{204} & \textbf{69} & \textbf{74.7\%} \\ \hline
  \end{tabular}
  \caption{דיוק השיוך המערכתי לפי מקצה.}
\end{table}

בטבלה זו לא בוצעה הפרדה כמותית בין סוגי הכשל. כלומר, מופע שלא נספר
כהצלחה יכול לנבוע מזיהוי רובוט חסר, איבוד עקבה, היעדר שיוך, או שיוך
לקבוצה שגויה. שגיאות אלו מפורטות בהמשך הפרק. מטרת המדד היא למדוד את
הצלחת המערכת מקצה לקצה, ולא לבודד את מקור השגיאה בכל רכיב בנפרד. בגרסה
עתידית ניתן לבצע ניתוח שגיאות מפורט יותר שיחלק את הכשלים לקטגוריות
כגון זיהוי רובוטים, מעקב, זיהוי ספרות ושיוך סופי.

משמעותו הפרקטית של נתון זה היא שבמדגם שנבדק, מתוך שישה רובוטים
הנמצאים בדרך כלל על המגרש, המערכת מצליחה לשייך נכון בממוצע כארבעה עד
חמישה רובוטים במצבים שבהם הם נראים על המסך וניתנים למעקב אנושי.

\begin{table}[H]
  \centering
  \begin{tabular}{|c|c|c|}
    \hline
    \textbf{רכיב} & \textbf{מדד} & \textbf{תוצאה} \\ \hline
    מודל זיהוי הרובוטים & \textenglish{mAP@50} & \textbf{0.97} \\ \hline
    מודל זיהוי הספרות &\textenglish{Precision}  ברמת הספרה & \textbf{0.75} \\ \hline
    שיוך סופי של עקבה לקבוצה & דיוק זיהוי ושיוך מתוך סרטון & \textbf{0.75} \\ \hline
    ביצועי הרצה & \textenglish{FPS} & \textbf{30--60} \\ \hline
  \end{tabular}
  \caption{סיכום התוצאות הכמותיות של המערכת.}
\end{table}

בנוסף למדדי הדיוק, המערכת מצליחה לייצר עבור כל קבוצה מסלול תנועה רציף
ומוחלק שאותו אפשר להציג לאחר המקצה לצורך ניתוח משחק.

מהתוצאות עולה שזיהוי הרובוטים הוא הרכיב היציב יותר, בעוד שזיהוי הספרות
הוא עדיין צוואר הבקבוק של המערכת. למרות זאת, השילוב של זיהויי ספרות
מצטברים לאורך זמן עם צבע הברית ורשימת הקבוצות האפשריות, מאפשר למערכת
להגיע לדיוק סופי גבוה משמעותית מזה שהיה מתאפשר מזיהוי ספרות בלבד.

\subsection{דוגמאות הצלחה}

בבדיקות המערכת נצפו כמה דפוסי הצלחה ברורים:

\begin{itemize}
\item \textbf{זיהוי ותחימה יציבים של ששת הרובוטים:} כאשר הרובוטים
  מרוחקים יחסית זה מזה ואינם חוסמים זה את זה, מודל הזיהוי והמעקב
  מצליחים לשמור על עקבות נקיות ויציבות לאורך זמן.

\item \textbf{התאוששות משגיאת ספרות בודדת:} גם כאשר פריים אחד מחזיר
  זיהוי חלקי או שגוי של הספרות, מנגנון ההצבעה המצטבר מונע מן המערכת
  לקפוץ מיד לזהות שגויה. לאחר כמה פריימים טובים יותר, המערכת מתכנסת
  שוב לקבוצה הנכונה.

\item \textbf{שיוך זהות נכון בזכות שילוב כמה רמזים:} במקרים רבים,
  הספרות לבדן אינן מספיקות. עם זאת, כאשר הן משולבות עם צבע הברית ועם
  רשימת שש הקבוצות האפשריות במקצה, אפשר להגיע לזהות נכונה גם אם רק חלק
  מן הספרות נקראו היטב.

\item \textbf{הפקת מסלולים שימושיים לניתוח:} גם במקרים שבהם היו חסרים
  פריימים בודדים או זיהויים לא מושלמים, האינטרפולציה וההחלקה הצליחו
  להפיק מסלול תנועה סביר, קריא ורציף לשימוש בתצוגת
  ה־\textenglish{Replay}.
\end{itemize}

דוגמאות וויזואליות להצלחות של זיהויי רובוטים וזיהויי ספרות מצורפות בנספח \ref{app:success} באיור \ref{fig:success}.

\subsection{דוגמאות כישלון וניתוח טעויות}

לצד ההצלחות, המערכת עדיין נכשלת בכמה סוגי מצבים שחוזרים על עצמם:

\paragraph{כשלים חזותיים:} הכשלים הבאים נובעים ממגבלות המידע החזותי
הזמין למודלים.

\begin{tabular}{|p{0.3\textwidth}|r|p{0.4\textwidth}|} \hline 
  \textbf{תמונה} & \textbf{סוג הכשל} & \textbf{השפעה על המערכת} \\ \hline
  \begin{minipage}[t]{\linewidth}
    \vspace{0pt}
    \includegraphics[width=\linewidth]{robots-failure-occluded}
    \vspace{0pt}
  \end{minipage}  & חסימת באמפר & הספרות או הרובוט כולו אינם נראים למשך זמן ממושך ולכן קשה לצבור ראיות מספיקות לזיהוי הקבוצה \\ \hline
  \begin{minipage}[t]{\linewidth}
    \vspace{0pt}
    \includegraphics[width=\linewidth]{digits-failure-blur}
    \vspace{0pt}
  \end{minipage} & טשטוש תנועה & זיהוי הספרות מחזיר לעיתים תוצאה חלקית או שגויה \\ \hline
  \begin{minipage}[t]{\linewidth}
    \vspace{0pt}
    \includegraphics[width=\linewidth]{digits-failure-gp}
    \vspace{0pt}
  \end{minipage} & שונות בין עונות & ביצועי המודלים עלולים לרדת כשהם
  מופעלים על עונה שלא הופיעה בנתוני האימון, בשל ההבדלים הגדולים במבנה
  הרובוטים, במגרש ובחפצי המשחק \\ \hline
\end{tabular}

\paragraph{כשלים מערכתיים:} הכשלים הבאים נובעים מההתנהגות וקבלת
ההחלטות של המערכת לאורך זמן.

\begin{tabular}{|p{0.15\textwidth}|p{0.4\textwidth}|p{0.3\textwidth}|} \hline
  \textbf{סוג הכשל} & \textbf{השפעה על המערכת} & \textbf{מנגנון התמודדות קיים} \\ \hline
  מספרי קבוצות דומים & עלול להיווצר בלבול בין שתי קבוצות מועמדות & דרישת מרווח מינימלי בין שתיהן \\ \hline
  החלפת \textenglish{Track IDs} & שיוך של רובוטים עלול להיפגע אחרי הצטלבות מסלולים או הסתרה ממושכת & איפוס השיוך וצבירת ראיות מחדש \\ \hline
\end{tabular}

דוגמאות הכישלון חשובות לא פחות מן ההצלחות, משום שהן מראות היכן המערכת
עדיין פגיעה ואילו שיפורים עתידיים יהיו בעלי התרומה הגדולה ביותר.

\subsection{מדריך למשתמש: תהליך השימוש והדגמת המערכת}

השימוש באפליקציית \textenglish{AutoScout} מתבצע בשלושה מסכים עיקריים:
מסך העלאת הסרטון, מסך התקדמות העיבוד, ומסך הצפייה בתוצאות.
מדריך זה מתאר את הפעולות שהמשתמש מבצע בכל אחד מהמסכים.

\subsubsection{מסך העלאת סרטון}

במסך הראשון מזינים את מספר המקצה ובוחרים את קובץ הווידאו שרוצים לנתח.
תחילה ממלאים את השדה \textenglish{Match Number}, ולאחר מכן לוחצים על
\textenglish{Select Video} כדי לבחור את הסרטון מהמחשב.

לאחר שנבחר קובץ וידאו, לוחצים על \textenglish{Upload} כדי לשלוח את הסרטון
לעיבוד. אם לא נבחר סרטון, כפתור ההעלאה אינו פעיל.

\begin{figure}[H]
  \centering
  \includegraphics[width=0.95\textwidth]{app-upload.png}
  \caption{מסך העלאת הסרטון והזנת מספר המקצה.}
\end{figure}

\subsubsection{מסך התקדמות העיבוד}

לאחר העלאת הסרטון מופיע מסך התקדמות העיבוד.
במסך זה ניתן לראות את מצב הניתוח של הסרטון בזמן שהשרת מעבד אותו.

המסך מציג את מספר הפריים הנוכחי מתוך כלל הפריימים, את קצב העיבוד
בפריימים לשנייה, ואת זמן הסיום המשוער. בנוסף מוצג פס התקדמות המאפשר
לעקוב אחר קצב ההתקדמות של הניתוח.

מסך זה נועד להראות שהמערכת עדיין עובדת ושאין צורך להעלות את הסרטון מחדש
בזמן שהעיבוד מתבצע.

\begin{figure}[H]
  \centering
  \includegraphics[width=0.95\textwidth]{app-progress.png}
  \caption{מסך התקדמות העיבוד של הסרטון.}
\end{figure}

\subsubsection{מסך צפייה בתוצאות}

לאחר סיום העיבוד מופיע מסך התוצאות.
במסך זה מוצג הסרטון לאחר הניתוח, יחד עם תוצאות המעקב והשיוך שהמערכת
הפיקה עבור הרובוטים.

במסך זה ניתן:
\begin{itemize}
    \item להפעיל ולעצור את הסרטון באמצעות כפתור הניגון.
    \item לעבור לנקודות זמן שונות במקצה באמצעות פס הזמן.
    \item לצפות בסימונים ובמספרי הקבוצות שהמערכת מציגה על גבי הסרטון.
    \item להשתמש בכפתור העליון כדי להציג או להסתיר את הסרטון המקורי מאחורי תוצאות המעקב.
    \item לחזור למסך הקודם באמצעות כפתור החזרה.
\end{itemize}

מסך זה מאפשר לבדוק את תוצאת המערכת בצורה חזותית. במקום להציג רק טבלה או
קובץ פלט, ניתן לראות את תוצאות המעקב והשייכות ישירות על גבי הסרטון.

\begin{figure}[H]
  \centering
  \includegraphics[width=0.95\textwidth]{app-replay.png}
  \caption{מסך צפייה בתוצאות לאחר סיום הניתוח.}
\end{figure}

\subsubsection{סיכום תהליך השימוש}

תהליך השימוש המלא במערכת הוא:
\begin{enumerate}
    \item מזינים מספר מקצה.
    \item בוחרים קובץ וידאו.
    \item לוחצים על \textenglish{Upload}.
    \item ממתינים לסיום העיבוד במסך ההתקדמות.
    \item צופים בתוצאות במסך הצפייה.
\end{enumerate}

\href{https://www.youtube.com/watch?v=x60OrcgcYnA}{\textenglish{AutoScout}: הדגמה} --- קישור לסרטון הדגמה קצר של המערכת בפעולה.

\subsection{סיכום פרק התוצאות}

התוצאות מראות שהמערכת חזקה בזיהוי הרובוטים, ושימושית גם בשיוך קבוצות
למרות חולשות קיימות בזיהוי הספרות.
